(1) 光从分束器到移动反射镜往返两次，镜位移 \(x\) 使光程差改变 \(2x\)，相位变化 \(\Delta\varphi=(2\pi/\lambda)2x=4\pi x/\lambda\)。等强理想干涉下某一输出口的光子份额为 \(p=(1+\cos\varphi)/2\)，另一个端口为 \(1-p\)。因此单端口光子速率 \(\Phi p\)，移动 \(\lambda/2\) 对应相位增加 \(2\pi\)。

(2) 积分时间 \(\tau\) 内探测数 \(N\) 的平均值 \(\bar N=\Phi\tau(1+\cos\varphi)/2\)，Poisson 方差 \(\operatorname{Var}N=\bar N\)。当输出位于最大或最小干涉强度时 \(\partial\bar N/\partial\varphi\) 为零，普通线性计数无法分辨小相位正负变化；在条纹斜率非零处，输出对镜位移的线性灵敏度才有意义。

(3) 由误差传播 \(\delta\varphi_{\rm one}=\sqrt{\bar N}/|\partial\bar N/\partial\varphi|=\sqrt{2(1+\cos\varphi)/(\Phi\tau)}/|\sin\varphi|\)。在 \(\varphi=\pi/2\) 处为 \(\sqrt{2/(\Phi\tau)}\)，对应 \(\delta x=\lambda\delta\varphi/(4\pi)\)。要注意在暗端极限必须重新考虑小计数分布与非线性估计器，不可单凭形式上 \(0/0\) 的误差传播公式宣称无限精度。

(4) 两个互补端口同时记录时 \(\bar N_1+\bar N_2=\Phi\tau\)，差分 \(D=N_1-N_2\) 的平均为 \(\Phi\tau\cos\varphi\)。独立泊松端口给 \(\operatorname{Var}D=\Phi\tau\)，因此在 \(\varphi=\pi/2\) 有 \(\delta\varphi_{\rm diff}=1/\sqrt{\Phi\tau}\)，比单一输出端口改善 \(\sqrt2\)。但总探测光子数有限，量子散粒噪声仍使相位分辨率按 \(1/\sqrt{\Phi\tau}\) 缩放，不能靠重复计算同一光子把噪声消除。
